\chapter{有限尺寸、无序与工艺误差鲁棒性}
\label{ch:robustness}

\section*{学习目标}
\begin{itemize}
  \item 理解为什么器件设计不能只优化名义值，以及“最优单点设计”和“鲁棒性设计”在目标函数上有何根本区别。
  \item 掌握孔半径无序、晶格常数漂移、刻蚀深度误差、侧壁倾角、表面粗糙、金属错位、电流孔径错位和阵列边缘截断如何进入 PCSEL 模型。
  \item 学会把几何误差、电学误差和热学误差统一组织成 tolerance analysis，并知道何时可用微扰观点，何时必须重做全波或多物理场计算。
\end{itemize}

\section*{先修知识}
建议已掌握 \cref{ch:fdtd-fem,ch:electrical,ch:eto-coupling,ch:single-mode-design} 中关于有限尺寸、注入非均匀性、热回写和大面积单模设计的内容。

\section*{本章路线图}
上一章已经把“大面积单模设计”写成了三层问题：单胞候选资格、有限孔阵横向可实现性、以及器件工作点下的真实胜负。现在必须进一步承认一个更现实的事实：真实器件从来不是名义 CAD 几何。孔半径会偏、晶格常数会漂、刻蚀会浅或过深、侧壁不是理想竖直、金属和电流孔径也未必和设计中心完全对齐。若书稿停留在名义值设计，那么它仍然更像一份研究草图，而不是一份器件教材。本章因此不把误差写成零散 checklist，而是从“误差如何进入模型”开始，随后讨论它们如何改变模式排序、辐射损耗、偏振、远场、注入均匀性和热分布，最后把这些物理变化组织成真正可执行的 tolerance analysis 与 robustness window 报告规范。

\section{为什么名义值最优不等于器件最优}
在很多纯光学优化问题里，设计者会先寻找一个名义参数点 $\bm{p}_0$，使某个性能指标在这个点上达到最好。例如，让被动辐射损耗最低、远场主瓣最窄、或者目标模与竞争模的频率分裂最大。这样的做法本身没有错，但它只回答了一个单点问题：若器件精确等于这一个参数点，它会怎样表现。真实器件设计要回答的是另一个问题：当参数在加工、公差、封装和工作过程中发生不可避免的偏移时，目标模式是否仍然赢得足够稳定。

于是，名义值最优和鲁棒性最优的目标函数天然不同。名义值优化追求
\begin{equation}
\max_{\bm{p}} \, F(\bm{p}),
\label{eq:nominal-opt}
\end{equation}
而鲁棒性设计更接近追求
\begin{equation}
\max_{\bm{p}} \, \mathbb{E}[F(\bm{p}+\delta \bm{p})]
\quad\text{或}\quad
\max_{\bm{p}} \, \min_{\delta \bm{p}\in \mathcal{U}} F(\bm{p}+\delta \bm{p}),
\label{eq:robust-opt}
\end{equation}
其中 $\delta \bm{p}$ 表示工艺与装配引入的偏差，$\mathcal{U}$ 是容差集合。\cref{eq:nominal-opt} 对应的是“找到理论最好点”；\cref{eq:robust-opt} 对应的是“找到在扰动下仍然可用的工作窗口”。二者可能重合，但在 PCSEL 中经常并不重合，因为目标模的优势常常建立在 Gamma 点简并分裂、有限阵列边缘控制、注入均匀性和热平衡这几个本来就敏感的机制之上。

\begin{IntuitionBox}
对 PCSEL 来说，最优单点设计关心的是“峰值有多高”，鲁棒性设计关心的是“平台有多宽”。器件导向研究更需要后者，因为加工出来的是一个有公差的结构，而不是一个数学点。
\end{IntuitionBox}

\section{误差如何进入模型：从参数向量到性能向量}
为了把各种误差放进同一框架，先把器件的名义参数写成
\begin{equation}
\bm{p}
=
\bigl(
r,\,
a,\,
h_{\mathrm{etch}},\,
\theta_{\mathrm{sw}},\,
\sigma_{\mathrm{rough}},\,
\Delta x_{\mathrm{metal}},\,
\Delta x_{\mathrm{ap}},\,
\bm{p}_{\mathrm{edge}},\,
\dots
\bigr),
\label{eq:param-vector}
\end{equation}
其中 $r$ 是孔半径，$a$ 是晶格常数，$h_{\mathrm{etch}}$ 是刻蚀深度，$\theta_{\mathrm{sw}}$ 是侧壁角，$\sigma_{\mathrm{rough}}$ 是表面粗糙度代理量，$\Delta x_{\mathrm{metal}}$ 表示金属图形相对孔阵的错位，$\Delta x_{\mathrm{ap}}$ 表示电流孔径相对目标模式中心的错位，$\bm{p}_{\mathrm{edge}}$ 描述阵列边缘截断方式。

与之对应，把关心的器件输出组织成性能向量
\begin{equation}
\bm{y}
=
\bigl(
\omega_m,\,
\alpha_{\mathrm{rad},m},\,
\Delta G_{\star j},\,
\Pi,\,
\mathcal{F}_{\mathrm{FF}},\,
\eta_{\mathrm{inj}},\,
T_{\max},\,
\dots
\bigr),
\label{eq:output-vector}
\end{equation}
其中 $\omega_m$ 是模式频率，$\alpha_{\mathrm{rad},m}$ 是辐射损耗，$\Delta G_{\star j}$ 是上一章定义的目标模相对竞争模的净裕量，$\Pi$ 可用来代表偏振纯度，$\mathcal{F}_{\mathrm{FF}}$ 代表远场特征量，$\eta_{\mathrm{inj}}$ 表示注入效率，$T_{\max}$ 表示峰值温升。鲁棒性问题的核心不是盯住某一个标量，而是研究
\begin{equation}
\delta \bm{p}
\longrightarrow
\delta \bm{y}
\label{eq:delta-map}
\end{equation}
这条映射。

若扰动足够小、模式身份不变且开放损耗变化不太强烈，可以先用一阶灵敏度近似
\begin{equation}
\delta y_i
\approx
\sum_j
\left.
\frac{\partial y_i}{\partial p_j}
\right|_{\bm{p}_0}
\delta p_j.
\label{eq:first-order-sensitivity}
\end{equation}
这一步的意义不是用线性模型替代真实器件，而是先快速识别“谁最敏感”。一旦出现模式排序改变、远场拓扑变化、注入路径重组或热热点迁移，\cref{eq:first-order-sensitivity} 往往就不再够用，必须回到全波或多物理场重算。

\begin{RigorousBox}
\cref{eq:first-order-sensitivity} 是工程上的局部近似。它成立的前提至少包括：扰动幅度小、模式标签可连续追踪、目标量对参数是局部光滑的。对强辐射开放模、模式简并附近、以及电热反馈导致的阈值翻转问题，这个近似必须谨慎使用。
\end{RigorousBox}

\section{几何无序的第一组：孔半径无序与晶格常数漂移}
孔半径无序（hole radius disorder）是最典型、也最容易在早期模型中被忽略的误差。对单胞来说，孔半径直接改变介电常数的平面傅里叶分量，因此会同时影响 Gamma 点附近耦合强度、频率位置和辐射通道强弱。若无序是全局一致的，例如整片晶圆上孔半径系统性偏小，那么它更像参数漂移；若无序是从孔到孔随机变化，那么它会打破原本依赖周期性的相干反馈，并可能把原本受对称性保护的模混合起来。

对有限器件来说，孔半径无序最先改变的通常不是“有没有模”，而是模式排序和远场纯净度。因为目标模与竞争模对单胞微扰的敏感方向并不相同，有些模更依赖某个特定 Fourier 分量维持低辐射与高相干，一旦这些分量在横向上随机起伏，目标模的辐射损耗就可能增加得更快。于是本来在名义值设计中占优的模式，可能在有限阵列里被别的模反超。

晶格常数变化（lattice constant variation）看起来只是比例尺变化，实际上影响更深。全局的 $a$ 漂移会整体移动归一化频率与辐射条件，从而改变工作波长相对于增益峰位和 light cone 的位置；局部的 $a(\rvec)$ 起伏则会造成 Bloch 相位积累失配，相当于在面内反馈网络中引入相位噪声。对大面积 PCSEL 来说，这种相位失配尤其关键，因为它会先破坏大面积相干，再表现为远场主瓣变宽、旁瓣上升或偏振不稳定。

从建模层级看，小幅全局 $r$ 或 $a$ 漂移通常可先用 PWE、CWT 或单胞 RCWA 做局部灵敏度；但只要问题变成孔到孔随机无序，或者设计本身工作在简并/近简并点附近，就应尽快转向有限尺寸全波模型来检查模排序和远场是否保持。若同时还关心阈值电流变化，则必须把这一步与后续的注入和热回写结合。

\section{几何无序的第二组：刻蚀深度、侧壁倾角与表面粗糙}
刻蚀深度误差（etch depth error）之所以危险，是因为它不只是改变平面内散射强度，还直接改变纵向层栈中的有效折射率剖面。换句话说，它同时碰了 \cref{ch:epitaxial-optics,ch:feedback} 的两条主线：既动了纵向模背景，也动了面内耦合与垂直辐射机制。若设计恰好依赖浅刻蚀或部分刻蚀来平衡辐射与导模约束，那么刻蚀深度误差对辐射损耗、上下辐射不对称和偏振都可能十分敏感。

侧壁角或锥度（sidewall angle / taper）进一步说明真实工艺不是理想二维截面。理想模型通常默认孔壁竖直，而真实刻蚀常导致上宽下窄或上窄下宽的三维轮廓。这会改变孔洞在不同高度处的有效横向尺寸，从而影响纵向不同场分量看到的等效散射强度。对近似 TE/TM 分类已经不够严格的 slab PCSEL 而言，这种三维几何变化还可能增强矢量模混合，使原本清晰的偏振选择变得更脆弱。

表面粗糙（surface roughness）则更像分布式弱散射源。若粗糙幅度较小、相关长度远小于模式横向尺度，工程上常把它先当成附加散射损耗或附加吸收代理，并研究其对 $Q$、阈值增益和线宽趋势的影响。但这只是第一层近似。因为粗糙不只改变总损耗，还可能打散原本规则的远场相位关系。对于大面积单模器件，这种影响有时会首先表现为远场旁瓣抬升、偏振纯度下降和批次间重复性变差。

这一组误差的共同特点是：它们都带有明显的三维性和开放辐射敏感性。因此，若只是想筛查灵敏度大小规律，可先在单胞 RCWA 或周期频域模型中扫 $h_{\mathrm{etch}}$ 与 $\theta_{\mathrm{sw}}$；若要判断远场、偏振和有限器件模排序是否真正稳健，通常必须进入三维有限阵列 FDTD 或有限元频域模型。表面粗糙若只作为附加损耗代理使用，应明确写出这是简化模型，而不是对真实散射过程的完整求解。

\begin{ExampleBox}
设某设计名义上依赖浅刻蚀来实现“目标模足够辐射、竞争模仍受抑制”的平衡。若刻蚀实际偏浅，目标模可能重新被过度束缚，导致出光下降；若刻蚀偏深，竞争模和目标模都可能打开更多辐射通道，净排序反而变差。这里最关键的不是总损耗单调上升还是下降，而是目标模相对于竞争模的损耗排序是否翻转。
\end{ExampleBox}

\section{装配与器件误差：金属错位、电流孔径错位和边缘截断}
到了器件层面，误差不再只发生在光子晶体孔阵本身。金属图形相对孔阵的错位（metal misalignment）会改变两个看似分离、实际上紧密耦合的问题。第一，它改变金属邻近吸收和上下边界的局部光学条件；第二，它改变电流进入有源区的路径和局部串联电阻。于是，同样的金属错位既可能改写辐射损耗，也可能改写注入均匀性和温升分布。

电流孔径错位（current aperture misalignment）更直接地触碰 \cref{ch:electrical} 建立的器件主线。对大面积 PCSEL 来说，目标模能否保持优势，很大程度上取决于高注入区域是否与目标模的主要场强区域重合。若孔径中心偏离目标模中心，哪怕冷腔模场本身十分漂亮，也可能因为局部模增益分布失配而让竞争模在边缘或局部热点处先获利。这正说明了为什么“光学模场分布合理”并不等于“器件注入均匀”。

阵列边缘截断方式（array edge truncation）则是有限尺寸问题最具体的落地形式。很多名义值设计默认孔阵是一个完美的矩形或圆形窗口，但实际加工中边缘往往会因为版图裁切、刻蚀邻近效应和保护环设置而形成不同的终止单元。边缘是有限阵列泄漏、散射和模式包络反射的交汇处，因此截断方式不只是版图收尾，而是有限器件边界条件的一部分。不同边缘截断可能导致完全不同的侧向泄漏排序，进而改变大面积单模性。

这一组误差几乎总是要求比单胞更高层级的模型。金属错位和孔径错位若只想看局部一阶趋势，可以先用简化的横向注入代理或损耗代理做 screening；但若要判断阈值电流、热点位置、模式竞争和远场稳定性，则必须进入电热耦合后的有限尺寸模型。边缘截断方式则天然属于有限器件全波问题，几乎不能仅靠无限周期模型回答。

\section{误差为什么会改变模排序、辐射、偏振、远场、注入与热}
到这里可以把各种误差背后的共同机制总结出来。第一类机制是\textbf{相位网络变化}。孔半径无序、晶格常数起伏和边缘截断会改变面内反馈网络的相位匹配，因此它们首先威胁的是模式排序和远场相干。第二类机制是\textbf{辐射通道重分配}。刻蚀深度、侧壁角和表面粗糙会改变目标模和竞争模向上、向下以及横向泄漏的比例，因此它们首先威胁的是辐射损耗、偏振和远场纯净度。第三类机制是\textbf{工作点重定位}。金属错位和电流孔径错位会改变载流子和温度的空间分布，因此它们首先威胁的是注入均匀性、热点位置和阈值电流。

这三类机制在真实器件里往往不是分开的。例如，金属错位先造成电流分布偏斜，进而产生温升不对称；温升不对称再通过折射率和增益谱漂移改变模式排序；模式排序变化又进一步改变近场和远场。也正因为这条链条存在，鲁棒性问题若只停留在冷腔名义值，就会系统性低估工艺误差的真实后果。

为了在设计中跟踪这些变化，一个实用写法是继续使用上一章的模裕量
\begin{equation}
\Delta G_{\star j}(\bm{p})
=
\bigl[g_{\mathrm{mod},\star}(\bm{p})-\alpha_{\star}(\bm{p})\bigr]
-
\bigl[g_{\mathrm{mod},j}(\bm{p})-\alpha_{j}(\bm{p})\bigr].
\label{eq:robust-margin}
\end{equation}
鲁棒性分析的核心不是只问 $\Delta G_{\star j}(\bm{p}_0)$ 是否为正，而是问在容差集合内
\begin{equation}
\Delta G_{\star j}(\bm{p}_0+\delta\bm{p}) > 0
\end{equation}
是否仍对足够多的扰动样本成立。若答案是否定的，那么名义设计再漂亮，也不该被称为鲁棒设计。

\section{哪些误差可以先用微扰观点，哪些必须重算}
这一节最容易被写成表格清单，但真正重要的是建立判断标准。能否先用微扰观点，不取决于误差名称本身，而取决于它是否满足三个条件：第一，扰动幅度小到不会改变拓扑和模式身份；第二，目标输出对参数局部光滑；第三，误差主要作用于已经分离好的单一物理模块，而不是同时重组光学、电学和热学链。

满足这些条件时，适合先做微扰或局部灵敏度分析的情形通常包括：小幅全局孔半径偏移、小幅全局晶格常数偏移、浅层刻蚀深度偏差、以及可被等效为附加损耗的弱粗糙。这些问题可先用 PWE/CWT/RCWA 或局部频域模型给出敏感方向，再挑最敏感的样本进入全波验证。

一旦出现下列情形，就应直接升级到全波或多物理场重算：孔到孔随机无序导致模式局域化趋势；刻蚀深度与侧壁角同时变化并显著改变上下辐射；金属错位和孔径错位导致注入路径重组；边缘截断方式改变；温升和注入反馈已足以改变模式排序。在这些问题里，局部一阶导数往往不再是决定性信息，因为真实器件响应可能由模翻转、热点迁移或远场拓扑变化主导。

\begin{table}[htbp]
  \centering
  \footnotesize
  \caption{常见误差的首选分析层级与最终验证层级。}
  \label{tab:robustness-model-choice}
  \setlength{\tabcolsep}{4pt}
  \begin{tabularx}{\textwidth}{@{}p{2.3cm}p{3.7cm}Y Y@{}}
    \toprule
    误差类型 & 首先关注的物理量 & 初筛模型 & 最终验证模型 \\
    \midrule
    孔半径全局偏移 & 频率漂移、辐射排序、偏振趋势 & PWE / CWT / 单胞 RCWA & 有限阵列全波 \\
    孔到孔随机半径无序 & 模排序、远场纯度、局域散射 & 有限阵列简化扫描 & 有限阵列全波，多样本统计 \\
    晶格常数全局漂移 & 归一化频率、增益峰失配 & PWE / 单胞频域 & 单胞开放模型加有限阵列抽查 \\
    刻蚀深度与侧壁角 & 辐射通道、上下不对称、偏振 & 单胞 RCWA / 频域 FEM & 三维有限阵列全波 \\
    表面粗糙 & 附加散射损耗、远场退化 & 损耗代理或局部频域 & 需要时做校准后的全波抽查 \\
    金属错位 & 金属吸收、注入路径、热点 & 简化损耗/注入代理 & 电-热-光耦合有限器件模型 \\
    电流孔径错位 & 模增益分布、竞争模获益 & 注入代理 + 冷腔模场叠加 & 漂移-扩散 + 热回写 + 光学重算 \\
    阵列边缘截断 & 边缘泄漏、横向包络、远场 & 简化边缘参数扫 & 有限阵列全波 \\
    \bottomrule
  \end{tabularx}
\end{table}

\section{如何做 tolerance analysis：不是盲扫，而是分层抽样}
真正可用的 tolerance analysis 很少是“所有参数都从最小到最大全组合扫一遍”，因为 PCSEL 的参数维度高、每次全波或多物理场求解成本也高。更有效的流程通常分四步。

第一步，先做参数分层。把误差按物理来源分成几何、装配、电学和热学四组，并区分系统误差与随机误差。系统误差更适合做 corner analysis，随机误差更适合做 Monte Carlo 或拉丁超立方采样。

第二步，先做局部灵敏度筛查。通过\cref{eq:first-order-sensitivity} 或有限差分灵敏度，识别对 $\Delta G_{\star j}$、远场主瓣、偏振纯度和峰值温升最敏感的少数参数。没有必要把所有低敏感度参数都立即带入最昂贵模型。

第三步，对高敏感参数做分层验证。若某参数主要影响单胞辐射，可先在 RCWA 或周期频域里扫；若某参数主要影响有限尺寸边缘泄漏，应直接在有限阵列全波里扫；若某参数影响注入和温升，则必须进入 CHARGE/HEAT 或 Semiconductor/Heat Transfer 一类模型。

第四步，才做联合统计。此时不是所有参数都等权，而是对关键参数赋予来自工艺的统计区间，构造样本集 $\{\bm{p}^{(s)}\}_{s=1}^{N_s}$，计算通过率
\begin{equation}
P_{\mathrm{pass}}
=
\frac{1}{N_s}
\sum_{s=1}^{N_s}
\mathbf{1}\!\left[
\Delta G_{\star j}^{(s)} > 0,\,
\Pi^{(s)}>\Pi_{\min},\,
\mathcal{F}_{\mathrm{FF}}^{(s)}\in \mathcal{A}_{\mathrm{FF}},\,
T_{\max}^{(s)}<T_{\max}^{\mathrm{allow}}
\right],
\label{eq:pass-rate}
\end{equation}
其中 $\mathbf{1}[\cdot]$ 是指示函数。\cref{eq:pass-rate} 的意义很直接：不是问“名义器件好不好”，而是问“在一批可能加工出来的器件里，有多大比例仍满足设计要求”。

若采用 Lumerical 路径，通常会用脚本串联 MODE/RCWA/FDTD 与 CHARGE/HEAT，对样本进行分阶段求解；若采用 COMSOL 路径，则更适合在统一几何参数化下做 parametric sweep，并在必要时把最贵的多物理场样本数量压缩到关键角点和代表性随机样本。

\begin{table}[htbp]
  \centering
  \footnotesize
  \caption{器件导向 PCSEL 设计中可直接起用的工艺容差定量判据。它们不是普适常数，而是“超出此量级就必须高度警惕模式身份和器件窗口被改写”的经验起点。}
  \label{tab:robustness_tolerance_start}
  \begin{tabularx}{\textwidth}{@{}p{2.5cm}p{2.3cm}p{3.0cm}Y@{}}
    \toprule
    参数 & 起始容差窗口 & 首先监视的量 & 设计判据 \\ \midrule
    晶格常数 $a$ & $\pm 2\%$ & 频率偏移、增益峰失配 & 目标模排序不翻转，工作波长仍留在目标增益窗口内 \\ 
    孔半径 $r$ & $\pm 5\%$ & 辐射损耗、偏振趋势、远场旁瓣 & $Q$ 与远场主瓣保持在允许窗口内 \\ 
    刻蚀深度 $h_{\mathrm{etch}}$ & $\pm \SI{10}{\nano\meter}$ & 上下辐射比、辐射损耗、偏振 & 目标模仍保持最小光学阈值或最大模裕量 \\ 
    波导层厚度 & $\pm 5\%$ & 有效折射率、纵模位置、重叠因子 & 纵向模身份不翻转，$\Gamma_{\mathrm{opt}}$ 变化仍可接受 \\ 
    量子阱位置 & $\pm \SI{20}{\nano\meter}$ & $\Gamma_{\mathrm{QW}}$、阈值增益 & 阈值链变化不超过设计裕量 \\ 
    \bottomrule
  \end{tabularx}
\end{table}

若希望把“鲁棒”写成更硬的工程语言，至少应给出一条推荐设计裕量：在名义设计点下，目标模相对首个竞争模的净裕量建议满足 $\Delta G_{\star j}>\SI{5}{\per\centi\meter}$；在容差边界或最坏样本下，仍应保持 $\Delta G_{\star j}>\SI{1}{\per\centi\meter}$。只要这条裕量在统计样本中频繁翻到零以下，就不应再把该结构称为鲁棒设计。

\begin{table}[htbp]
  \centering
  \footnotesize
  \caption{鲁棒性报告时建议明确给出的量化判据。表中符号应由具体项目自行替换成明确门槛。}
  \label{tab:robustness_quantitative}
  \begin{tabularx}{\textwidth}{@{}p{2.8cm}p{3.1cm}Y@{}}
    \toprule
    判据类别 & 推荐量化写法 & 用途 \\ \midrule
    模身份保持 & 目标模在全部样本中与同一模族标签持续对应 & 防止把模式翻转误写成“参数敏感但可接受” \\
    模裕量下界 & $\min_s \Delta G_{\star j}^{(s)} > 0$ 或给出最坏样本下界 & 防止只汇报均值而不说明最差情形 \\
    远场窗口 & $\mathcal{F}_{\mathrm{FF}}^{(s)} \in \mathcal{A}_{\mathrm{FF}}$ 的通过率或最差偏离量 & 把远场要求从主观“看起来差不多”变成窗口判据 \\
    热安全边界 & $\max_s T_{\max}^{(s)} < T_{\max}^{\mathrm{allow}}$ & 防止只从光学视角宣布鲁棒 \\
    良率 / 通过率 & $P_{\mathrm{pass}}$ 或角点最差样本通过情况 & 把“鲁棒”写成明确概率或最坏情形语言 \\
    \bottomrule
  \end{tabularx}
\end{table}

\section{如何报告 robustness window}
鲁棒性报告最容易犯的错误，是只汇报“最优点”或者只给“某个参数可容忍 $\pm x$”而不说明其它参数是否同时固定。真正有用的 robustness window 至少应包含四层信息。

第一，名义设计点。必须给出 $\bm{p}_0$ 及其对应的名义性能。

第二，单参数窗口。说明在其它参数固定时，每个关键参数可接受的变化区间，例如孔半径、刻蚀深度、孔径错位等。

第三，联合窗口或通过率。说明当多个关键误差同时存在时，器件通过率是多少，或者最坏情况下的性能下界是多少。

第四，失效模式说明。明确器件首先因为什么失效，是模排序翻转、远场恶化、偏振崩坏、局部过热，还是注入均匀性丢失。没有失效模式说明的容差窗口，往往不利于后续工艺沟通。

在文风上，最推荐的报告方式不是说“该设计对工艺误差较鲁棒”，而是更具体地说：“在给定的孔半径、刻蚀深度与孔径错位统计区间下，目标模相对首个竞争模的净裕量仍保持为正，远场主瓣仍位于允许角窗内，峰值温升未超过设计上限，因此该方案在所定义的工艺窗口内通过率达到某一区间。”如果没有统计样本支持，就不应使用“鲁棒”这个词。

\begin{PitfallBox}
只给名义参数点，不给通过率、不说首先在哪种误差下失效，这样的“优化结果”通常只适合内部草图，不适合作为器件设计结论。鲁棒性设计的关键不是把均值做高，而是把失效边界说明白。
\end{PitfallBox}

\section{把鲁棒性设计接回工作流与失败模式}
现在可以把本章放回全书主线。\cref{ch:single-mode-design} 说明了为什么单模设计不能只盯名义值；本章进一步说明，即便名义值设计已经跨过单胞、有限阵列和器件工作点三层关口，仍然必须检查它对工艺误差是否稳健。\cref{ch:workflow} 中的工作流，则应在“名义结构通过电热光验证”之后显式加入 tolerance analysis 阶段，而不是把它留到最后当成可有可无的附加项。\cref{ch:failure-modes} 中的许多失败模式，其实正是在鲁棒性分析缺席时才暴露得最晚、代价也最高。

因此，器件导向的 PCSEL 设计流程更合理的顺序应是：先做名义值筛选，再做有限尺寸验证，再做器件级电热光回写，最后用 tolerance analysis 检查这个方案是否有足够宽的 robustness window。若一个方案只有名义值漂亮，却没有工艺窗口，它更适合作为物理演示结构，而不适合作为器件候选。

\section*{本章小结}
PCSEL 的鲁棒性问题不是简单的“公差补充说明”，而是器件设计的核心组成部分。孔半径无序、晶格常数漂移、刻蚀深度误差、侧壁角、表面粗糙、金属错位、电流孔径错位和边缘截断，都会通过面内相位网络、辐射通道重分配以及注入与温升重定位三条路径改变模式排序、辐射损耗、偏振、远场、注入均匀性和热分布。小扰动、单模块、模式身份稳定的问题可先用微扰观点筛查；一旦涉及模式翻转、有限尺寸边界、注入路径重组或热反馈，就必须升级到全波或多物理场重算。真正的 tolerance analysis 追求的不是名义值最优，而是在工艺扰动下仍有足够通过率的 robustness window。

\section*{练习题}
\begin{enumerate}
  \item \basicq 说明为什么孔半径无序与晶格常数随机起伏对大面积 PCSEL 的危害，往往首先表现为远场和模式排序问题，而不是单一频率漂移。

  \item \basicq 针对刻蚀深度误差，分别说明单胞 RCWA、有限阵列 FDTD/FEM 和电热耦合模型各自能回答什么、不能回答什么。

  \item \midq 给出一个你自己的 tolerance analysis 计划，至少包含参数分层、灵敏度筛查、联合抽样和通过率定义四步。

  \item \hardq 解释为什么“金属错位”与“电流孔径错位”不能仅当作几何误差处理，而应看作光学、电学和热学的耦合误差。
\end{enumerate}

\section*{建议继续阅读}
\begin{itemize}
  \item 下一章将先把“名义设计是否成功”推进到起振以上的动态稳定性、调制与噪声问题，再在随后的性能章节中把这些动态结论翻译成线宽、相干性、亮度和光束质量等可观测指标。 这条阅读的重点是把本章的输出顺着主线接到后续模块，而不是停成孤立知识点。

  \item 建议与 \cref{ch:workflow,ch:failure-modes} 对照阅读：前者告诉你 tolerance analysis 应放在工作流哪里，后者告诉你缺少鲁棒性分析时最容易出现哪些误判。 这条阅读的重点是把本章的输出顺着主线接到后续模块，而不是停成孤立知识点。
\end{itemize}


